
To enforce waste segregation, Local Government Units (LGUs) worldwide relied on a policy mix of economic incentives, regulatory penalties, and educational campaigns. The critical task for policymakers was to find the optimal combination and intensity of these levers—one that maximized public compliance and segregation rates without exceeding finite public budgets \citep{HeFu2021, Torkayesh2021}. This presented a complex optimization problem, as these policy levers were not only budget-dependent but also highly interconnected. This research, therefore, modeled the LGU's strategic choice across three such interconnected, budget-dependent policy levers.

\subsection{Economic and Regulatory Levers}
\setlength{\parindent}{2em}

Economic instruments, frequently framed as reward-penalty schemes, were among the most powerful direct drivers of compliance because they immediately altered the financial cost-benefit analysis of the household segregation decision \citep{MuZhang2021, Zhao2022}. While \citet{Cheng2022} focused on construction waste and \citet{Wang2023} analyzed closed-loop supply chains, both concluded that government-led incentive-punishment mechanisms were essential for rationalizing waste reduction behavior. The literature consistently confirmed that a hybrid approach—one combining both incentives and fines—was more effective than relying on either instrument alone \citep{Chen2023, MuZhang2021}. \citet{Rathore2021} further argued that identifying these suitable motivational mechanisms was critical for the success of any reverse logistics or collection system. This dual strategy effectively leveraged both the psychological gain associated with rewards and the powerful aversion to loss associated with penalties.

A critical complication, however, was that the effectiveness of these financial policies was not uniform. The impact of both incentives and penalties varied significantly based on household heterogeneity \citep{Chen2023, Zheng2022}. For instance, low-income households tended to be more sensitive to the disutility of fines, while the perceived benefit of an incentive might have been modulated by a household's education level or the perceived complexity of the program \citep{Zheng2022}. This finding directly justified the methodological necessity of an Agent-Based Model (ABM), which could model distinct household agents whose sensitivity to the LGU’s financial actions was weighted by their synthetic socio-economic profiles.

Finally, these economic policies had to be optimized not only for effectiveness but also for cost and robustness against uncertainty \citep{Gentile2022}. Robust optimization models in waste management explicitly sought to minimize the ``price of robustness''—that is, the extra cost incurred to protect the system against uncertain parameters, such as fluctuating waste volumes \citep{Gentile2022}. This concept directly mirrored the LGU’s practical constraint: the need to find a policy balance that avoided over-spending on enforcement or incentive schemes that, while effective, might yield diminishing returns \citep{Gentile2022}.

\subsection{Educational and Behavioral Levers}
\setlength{\parindent}{2em}

Educational strategies represented a critical, non-monetary intervention essential for achieving the long-term sustainability of Solid Waste Management (SWM) programs. Their fundamental value lay in their ability to address the root behavioral challenges that often undermined the success of technical or financial policies alone \citep{Moeini2023}. Complementing traditional education, \citet{LoanBalanay2023} advocated for the application of Nudge Theory, suggesting that subtle architectural changes and positive reinforcements could be as effective as strict mandates in reinforcing waste separation habits.

For a Local Government Unit (LGU), educational campaigns were the primary tool for dynamically influencing household behavior, particularly by targeting the core constructs of the Theory of Planned Behavior (TPB). These campaigns could cultivate a more positive Attitude toward segregation by clearly conveying its environmental importance. They also enhanced Perceived Behavioral Control (PBC) by providing specific, practical knowledge on how to segregate properly, thereby increasing residents' confidence that the action was feasible. Furthermore, education reinforced Subjective Norms (SN) by increasing social awareness and fostering a community-wide expectation of compliance \citep{Vorobeva2022}.

The impact of these educational efforts extended beyond mere awareness, playing a significant role in the adoption of new systems. Studies showed that ``soft'' behavioral factors, such as an established pro-environmental behavior (PEB) and a sense of empowerment, were crucial drivers for household participation. Notably, this influence persisted even when financial incentives were part of the policy mix \citep{Vorobeva2022}. This supported a modeling approach where investment in education—by enhancing these foundational behavioral factors—improved the general willingness of agents to participate in and comply with LGU programs, complementing other interventions.

An integrated Agent-Based Modeling (ABM) framework was particularly well-suited for simulating these complex policy interactions. ABM had been successfully validated in prior research for its ability to simulate community-level behavioral responses and visualize compliance patterns under multiple incentive policies \citep{Ma2023}. This established precedent provided confidence that the model could accurately capture and predict the complex, emergent results of a hybrid policy that combined both educational interventions and financial incentives.

\subsection{Resource Allocation Heuristics in Public Policy}
\setlength{\parindent}{2em}
The optimization of resource allocation in public policy is frequently bounded by established operational and economic heuristics, ensuring that interventions remain efficient and pragmatic under real-world constraints. Two foundational concepts in this domain are the Theory of Constraints (TOC) and the Law of Diminishing Marginal Returns.

Originating from operations research, the Theory of Constraints postulates that the overall performance of any complex system is strictly limited by its weakest link, or ``bottleneck'' \citep{Goldratt1984}. In the context of public administration and policy deployment, TOC translates to ``worst-first'' or triage allocation strategies. When an overarching system---such as municipal-wide environmental compliance---is failing, evenly distributing limited capital across all sectors yields suboptimal results. Instead, TOC dictates that resources must be overwhelmingly redirected to elevate the most critical bottleneck until the system stabilizes \citep{Umble2006}.

Conversely, the Law of Diminishing Marginal Returns governs the upper limits of policy effectiveness. In environmental economics, initial investments in compliance (e.g., capturing the first 60\% to 70\% of a population's adherence) yield high returns at a relatively low cost due to the presence of early adopters and accessible targets \citep{Tietenberg2018}. However, as compliance approaches absolute saturation, the cost to convert the remaining ``hardcore non-compliers'' increases exponentially. \citet{Rondinelli2000} note that traditional command-and-control environmental policies inevitably experience diminishing marginal returns, where extreme budget expenditures are required for statistically negligible behavioral shifts. Therefore, optimizing public policy requires dynamic thresholds: identifying the exact point where a population becomes saturated and reallocating the budget before marginal returns become negative.